%%%%%%%%%%%%%%%%%%%%%%%%%%%%%%%
\chapter{Higher-order pressure boundary condition}
\label{chap:hopressure-bndcond}

\definecolor{nekblue}{HTML}{24527A}
\definecolor{neklightblue}{HTML}{E9F2F8}
\definecolor{nekred}{HTML}{B44B4B}
\definecolor{neklightred}{HTML}{F8EBEB}
\definecolor{nekgreen}{HTML}{3E7654}
\definecolor{neklightgreen}{HTML}{EAF4ED}
\definecolor{nekgray}{HTML}{535B63}
\definecolor{neklightgray}{HTML}{F2F3F4}
\providecommand{\vect}[1]{\boldsymbol{#1}}

The higher-order pressure boundary-condition operator,
\texttt{HOPressureBndCondOp} of the incompressible Navier--Stokes solver,
constructs a consistent Neumann datum for the pressure Poisson equation in
the velocity-correction scheme. It evaluates the curl--curl of an
extrapolated velocity on the boundary-adjacent volume elements, traces the
resulting strong-form vector flux to boundary quadrature points, contracts it
with the outward normal, forms a boundary inner product, and scatters that
weak boundary load back into the pressure coefficient right-hand side. This
chapter describes the continuous condition, its time-discrete form, and each
discrete data movement implemented in \path{HOPressureBndCondOpImpl.hpp}, in
the directory \path{BoundaryConditions/HOPressureBndCond} of the solver.

\section{Purpose of the condition}

Write the incompressible momentum equation in the form
\begin{equation}
  \frac{\partial \vect{u}}{\partial t}
  = \vect{r} - \nabla p
    - \nu\,\nabla\times(\nabla\times\vect{u}),
  \label{hopbc:eq:momentum-rotational}
\end{equation}
where \(\vect{r}\) denotes the explicit strong-form contribution, normally
body forcing minus advection.  The rotational viscous form follows from
\begin{equation}
  \nabla^2\vect{u}
  = \nabla(\nabla\cdot\vect{u})
    - \nabla\times(\nabla\times\vect{u})
  \simeq -\nabla\times(\nabla\times\vect{u})
\end{equation}
for a divergence-free velocity field.

Taking the normal component of \eqref{hopbc:eq:momentum-rotational} gives the
consistent pressure Neumann condition
\begin{equation}
  \boxed{
  \frac{\partial p}{\partial n}
  = \vect{n}\cdot
    \left(
      \vect{r}
      - \nu\,\nabla\times(\nabla\times\vect{u})
      - \frac{\partial\vect{u}}{\partial t}
    \right).}
  \label{hopbc:eq:continuous-hop}
\end{equation}
It is called a higher-order pressure condition because it retains the
explicit, viscous, and acceleration terms dictated by the momentum equation;
it is not the commonly used approximation \(\partial p/\partial n=0\).

\section{The exact algebra evaluated by the code}

Let
\begin{center}
\begin{tabular}{@{}>{$}l<{$} p{0.68\linewidth}@{}}
\toprule
\text{Symbol} & \text{Implementation quantity} \\
\midrule
s & \code{scale}, set by \code{SetScale} \\
\nu & \code{kinvis}, set by \code{SetKinvis} \\
\vect{R} & \code{rhs}, the extrapolated explicit contribution passed to
\code{Update} \\
\vect{V} & \code{velocity}, the extrapolated velocity quantity used by the
curl--curl operator \\
\vect{P} & \code{predictor}, the complete implicit-stage right-hand side \\
\vect{u}_D(t) & prescribed velocity evaluated from the boundary expressions
at the target time, set by \code{SetTime} \\
\vect{C}(\vect{V}) & \(\nabla\times(\nabla\times\vect{V})\) \\
\bottomrule
\end{tabular}
\end{center}

\code{Update} evaluates
\begin{equation}
  g = s\,\vect{n}\cdot
      \left[\vect{R}-\nu\vect{C}(\vect{V})\right],
  \label{hopbc:eq:simple-overload}
\end{equation}
where \(g=\partial p/\partial n\).

\code{UpdateWithAcceleration}, which \code{VelocityCorrectionScheme} calls,
adds the acceleration to the flux of \code{Update}:
\begin{align}
  \vect{F}_{\mathrm{base}}
    &=s\left[\vect{R}-\nu\vect{C}(\vect{V})\right],\\
  \vect{F}_{\mathrm{acc}}
    &=s\left[\vect{P}-\vect{R}\right].
\end{align}
Their explicit contributions cancel exactly, so it takes the predictor
\(\vect P\) and not \(\vect R\). It then subtracts the prescribed boundary
velocity, producing
\begin{equation}
  \boxed{
  g=s\,\vect{n}\cdot
    \left[\vect{P}-\vect{u}_D(t)-\nu\vect{C}(\vect{V})\right].}
  \label{hopbc:eq:implemented-hop}
\end{equation}

A velocity component \(d\) that is not Dirichlet on a region, such as a
Neumann component of a slip wall or of an inflow, has no prescribed boundary
velocity and takes no acceleration there: its term in the flux is
\(s\,n_d\left[R_d-\nu C_d(\vect{V})\right]\). On such a region the operator
takes \(u_{D,d}=0\) and subtracts \(s\,n_d\left[P_d-R_d\right]\) from
\eqref{hopbc:eq:implemented-hop}, so \code{UpdateWithAcceleration} also reads
\(\vect R\). \code{VelocityCorrectionScheme} keeps the history of the
explicit forcing from which \(\vect R\) is extrapolated only on a rank
holding such a region, as \code{HasNonDirichletVelocity} reports.

In the current solver call,
\begin{equation}
  s=\frac{1}{\gamma\Delta t},
\end{equation}
with \(\gamma\Delta t\) the implicit step size of the stage.  The predictor
contains time-integrator-scaled explicit terms.  \code{VelocityCorrectionScheme}
keeps histories of the velocity and of the explicit forcing, and combines them
into \(\vect{V}\) and \(\vect{R}\) with the coefficients \(c_k\) of the explicit
contributions of the stage, scaled to sum to \(\gamma\Delta t\):
\begin{equation}
  \vect{V}=\frac{\gamma\Delta t}{\sum_k c_k}\sum_k c_k\,\vect{u}^{(k)},\qquad
  \vect{R}=\frac{\gamma\Delta t}{\sum_k c_k}\sum_k c_k\,\vect{r}^{(k)}.
\end{equation}
For a multistep IMEX scheme the \(c_k\) are \(\gamma\) times the
extrapolation weights and sum to \(\gamma\), so the scaling is \(\Delta t\).
For a later stage of an IMEXdirk scheme, such as the startup of a multistep
scheme, they are quadrature weights whose sum is not \(\gamma\); without the
scaling, \(s\vect{V}\) would not reduce to the velocity, and even a steady
solution would get a wrong flux.  Accordingly,
\(s(\vect{P}-\vect{u}_D)\) supplies the discrete negative acceleration in
\eqref{hopbc:eq:continuous-hop}, while
\(-s\nu\vect{C}(\vect{V})\) supplies the consistently extrapolated rotational
viscous term.  Keeping \eqref{hopbc:eq:implemented-hop} in terms of the actual stage
quantities is important: \(\vect{P}\), \(\vect{R}\), and \(\vect{V}\) are not
all unscaled instantaneous physical fields.

\section{Geometric setup}

Only pressure regions explicitly opted in via a Neumann boundary condition
tagged \code{USERDEFINEDTYPE="H"} are selected for higher-order treatment,
the tag the IncNavierStokesSolver also uses. An ordinary, untagged Neumann
pressure region keeps its plain session-prescribed value instead. The
operator is added to the pressure system with
\code{AddNeumannBCOp(op, "H")}, so the default Neumann operator of the
system leaves the tagged regions out and their session \code{VALUE} is not
applied. For every
selected region, \texttt{SetVelocityContext} matches each element of the
boundary expansion to the velocity volume element owning it as a trace.
Global geometry identifiers are used to record
\begin{equation*}
  \text{boundary face}
  \longleftrightarrow
  (\text{owning volume element }K,\ \text{local face }f).
\end{equation*}
A restricted volume expansion list is built from the unique set of elements
touching such faces.  Consequently, the expensive derivative operation is
performed only on boundary-adjacent elements.

\begin{figure}[htbp]
\centering
\begin{tikzpicture}[
  >=Latex,
  volpt/.style={circle,fill=nekblue,inner sep=1.55pt},
  bndpt/.style={circle,fill=nekred,inner sep=1.85pt},
  note/.style={font=\small,align=center}
]
  \coordinate (a) at (0,0);
  \coordinate (b) at (5.0,0.35);
  \coordinate (c) at (4.45,3.7);
  \coordinate (d) at (0.55,3.15);
  \fill[neklightblue] (a)--(b)--(c)--(d)--cycle;
  \draw[very thick,nekblue] (a)--(b)--(c)--(d)--cycle;
  \draw[line width=3.2pt,nekred] (b)--(c);

  \foreach \x/\y in {1.0/0.85,2.35/0.95,3.65/1.08,
                       1.1/1.85,2.35/2.05,3.55/2.25,
                       1.2/2.72,2.35/2.91,3.45/3.05}
    \node[volpt] at (\x,\y) {};
  \foreach \x/\y in {4.90/0.78,4.78/1.55,4.66/2.32,4.54/3.08}
    \node[bndpt] at (\x,\y) {};

  \node[note,nekblue] at (2.25,1.55)
    {volume quadrature values\\on the owning element $K$};
  \node[note,nekred,right=5mm] at (5.0,2.1)
    {pressure boundary\\quadrature points on $f$};
  \draw[->,thick,nekred] (3.6,1.65) to[bend left=18]
    node[above,sloped,font=\small] {boundary trace extraction} (4.64,2.30);
  \draw[->,thick,nekgreen] (5.2,2.85)--++(0.95,0.3)
    node[right,note,nekgreen] {outward normal\\$\vect n$};
  \node[below=5mm of a,anchor=west,note]
    {Only elements adjacent to higher-order pressure faces are gathered.};
\end{tikzpicture}
\caption{The curl--curl derivative is evaluated on the volume element.  A
shared boundary-trace-extraction operator then moves the resulting vector
quantity to the boundary quadrature grid.}
\label{hopbc:fig:geometry}
\end{figure}

\section{Volume-to-boundary evaluation}

\subsection{Gathering boundary-adjacent volume elements}

At each update, the velocity data on the selected elements is copied from the
full solver field into the restricted field:
\begin{equation}
  \vect{V}_{\mathrm{restricted}}|_K
  \leftarrow \vect{V}_{\mathrm{full}}|_K.
  \label{hopbc:eq:volume-gather}
\end{equation}
This is an indexed copy of every quadrature point and velocity component; it
does not project or interpolate the values.  The gather understands the
interleaved element layout used by Serial, AVX, CUDA, and HIP execution.

\subsection{Curl--curl on the restricted volume}

The element operator computes
\begin{equation}
  \vect{C}_K
  =\nabla\times(\nabla\times\vect{V}_K)
  \label{hopbc:eq:curlcurl}
\end{equation}
at volume quadrature points.  This operation must occur before tracing to the
boundary: its derivatives require the volume polynomial, not merely its face
values.

\subsection{Trace interpolation and normal contraction}
\label{hopbc:sec:trace-and-normal}

Each region is traced by its own instance of the library helper
\texttt{PhysTraceExtractBndCond} in \texttt{MultiRegions::detail}
(\S\ref{sec:phystraceextractbndcond}), which extracts a
volume field onto the physical points of a continuous boundary region,
including regions of mixed shape and order and the Gauss--Lobatto points of
collapsed directions. Its volume expansion list is the restricted list of
boundary-adjacent elements, which is built from an expansion-information map
and has no mesh graph, so the mesh graph of the velocity field is passed to
it. Writing \(\mathrm{Extract}_f(\cdot)\) for this
operator's action on face \(f\),
\begin{equation}
  v_f(x_a)=\mathrm{Extract}_f\bigl(v_K\bigr)(x_a),
  \label{hopbc:eq:trace-extract}
\end{equation}
where \(a\) indexes the boundary region's own quadrature points.

Because extraction is linear, the base and acceleration contributions are
combined \emph{before} tracing rather than after: one volume-domain kernel
first forms the combined flux \(s\,(\vect P-\nu\vect C)\) (or
\(s\,(\vect R-\nu\vect C)\) for \code{Update}) on the
boundary-adjacent volume elements, since
\(\mathrm{Extract}_f(\vect P)-\nu\,\mathrm{Extract}_f(\vect C)
=\mathrm{Extract}_f(\vect P-\nu\vect C)\). A single
\code{PhysTraceExtractBndCond::Apply} call per region then traces that
combined flux to boundary points, and a separate kernel contracts it with the
outward normal:
\begin{equation}
  g^f_a
  =s\,\vect n^f_a\cdot
   \mathrm{Extract}_f\bigl(\vect P-\nu\vect C\bigr)(x_a)
   -s\,\vect n^f_a\cdot\vect u_D(x_a,t).
  \label{hopbc:eq:pointwise-kernel}
\end{equation}
The prescribed velocity \(\vect u_D(t)\) is evaluated directly on the boundary
expansion (never traced from the volume), and the outward normals are stored
on the same boundary quadrature grid.

\section{Boundary-to-volume weak assembly}

The scalar array \(g^f_a\) is still a physical-space boundary field.  For
every boundary expansion of positive dimension, the boundary
\code{IProductWRTBaseOp} computes
\begin{equation}
  b^f_j
  =\int_{\Gamma_f}\phi^f_j\,g^f\,\mathrm dS
  \simeq
   \sum_a w_a J^f_a\,\phi^f_j(x_a)g^f_a.
  \label{hopbc:eq:boundary-iproduct}
\end{equation}
This is an inner product, not a forward transform: \(b^f_j\) is a weak-load
coefficient.  For a zero-dimensional boundary, such as an endpoint of a
one-dimensional domain, the physical value is copied directly.

The per-region boundary coefficient arrays are concatenated and rearranged
into the execution-block layout of the pressure field.  Finally, the
boundary-to-element maps scatter-add them into the pressure coefficient RHS:
\begin{equation}
  \boxed{
  \vect r_{\mathrm{pressure}}
  \mathrel{+}=
  \sum_{f\subset\Gamma_N}
  E_f^{\mathsf T}B_f^{\mathsf T}W_fJ_f\,\vect g_f.}
  \label{hopbc:eq:weak-scatter}
\end{equation}
Here \(B_f\) is the boundary basis evaluation matrix and \(E_f\) embeds face
coefficients into the owning volume expansion.  Orientation signs are applied
where required.  When two boundary pieces target the same volume coefficient,
for example at a corner shared by two Neumann faces, their contributions are
grouped, summed, and written once.

Thus, ``back to volume'' does \emph{not} mean constructing another physical
volume field.  It means embedding the boundary weak load into the coefficient
right-hand side consumed by the pressure linear system.

\begin{figure}[htbp]
\centering
\begin{tikzpicture}[
  >=Latex,
  node distance=5mm and 7mm,
  box/.style={draw=nekblue,very thick,rounded corners=2pt,
              fill=neklightblue,minimum height=11mm,text width=31mm,
              align=center,font=\small},
  redbox/.style={box,draw=nekred,fill=neklightred},
  greenbox/.style={box,draw=nekgreen,fill=neklightgreen},
  arrow/.style={->,very thick,nekgray},
  label/.style={font=\scriptsize,align=center,fill=white,inner sep=1pt}
]
  \node[box] (full) {full volume velocity\\$\vect V$};
  \node[box,right=of full] (gather) {gather selected\\volume elements};
  \node[box,right=of gather] (curl) {volume operator\\$\nabla\times\nabla\times\vect V$};
  \node[redbox,below=11mm of curl] (trace) {trace interpolation\\and flux fusion};
  \node[redbox,left=of trace] (normal) {normal flux at\\boundary points $\vect g_f$};
  \node[redbox,left=of normal] (iprod) {boundary inner product\\$B_f^{\mathsf T}W_fJ_f\vect g_f$};
  \node[greenbox,below=11mm of iprod] (compact) {compact and gather\\boundary coefficients};
  \node[greenbox,right=of compact] (scatter) {signed trace-to-element\\scatter $E_f^{\mathsf T}$};
  \node[greenbox,right=of scatter] (rhs) {pressure coefficient\\right-hand side};

  \draw[arrow] (full)--(gather);
  \draw[arrow] (gather)--(curl);
  \draw[arrow] (curl)--(trace);
  \draw[arrow] (trace)--(normal);
  \draw[arrow] (normal)--(iprod);
  \draw[arrow] (iprod)--(compact);
  \draw[arrow] (compact)--(scatter);
  \draw[arrow] (scatter)--(rhs);

  \draw[->,thick,nekred] ($(trace.north)+(0,5mm)$)--(trace.north)
    node[midway,right,label] {$\vect P$ or $\vect R$ from the owning volume};
  \draw[->,thick,nekred] ($(normal.north)+(0,5mm)$)--(normal.north)
    node[midway,right,label] {$\vect n,\ \vect u_D(t)$ on the boundary};
\end{tikzpicture}
\caption{Runtime data flow. Blue operations use volume data, red operations
produce the physical and weak boundary flux, and green operations return the
load to the pressure coefficient system.}
\label{hopbc:fig:dataflow}
\end{figure}

\section{Relation to the pressure Poisson weak form}

The redesign pressure solver uses a positive stiffness operator and applies a
factor of \(-1\) to the volume inner product.  For the strong equation
\begin{equation}
  \nabla^2p=f_p,
\end{equation}
integration by parts gives
\begin{equation}
  (\nabla p,\nabla v)_\Omega
  =-(f_p,v)_\Omega
   +\left\langle\frac{\partial p}{\partial n},v\right\rangle_{\Gamma_N}.
  \label{hopbc:eq:poisson-weak}
\end{equation}
The higher-order operator supplies precisely the positive boundary term in
\eqref{hopbc:eq:poisson-weak}.  Its \code{Apply} method therefore adds to, rather
than overwrites, the coefficient RHS assembled from the volume forcing.

\section{Implementation map}

\begin{center}
\small
\begin{tabular}{@{}p{0.37\linewidth}p{0.59\linewidth}@{}}
\toprule
\textbf{Implementation routine} & \textbf{Mathematical or data-layout role} \\
\midrule
\code{BuildBoundaryData} & Select pressure regions opted in via
\code{USERDEFINEDTYPE="H"}; construct boundary physical/coefficient fields
and continuous-Galerkin scatter maps. \\
\code{v\_SetVelocityContext} & Build the velocity-dependent data below, once,
before time stepping. \\
\code{BuildBoundaryElementData} & Match faces to owning elements; create the
restricted volume list, volume gather maps, the \code{CurlCurlOp}, and one
\code{PhysTraceExtractBndCond} per region. \\
\code{BuildBoundaryVelocityData} & Store the outward normals and set up the
\code{ExpressionOp} of the prescribed boundary velocity for every region;
on a region with a velocity component that is not Dirichlet, also store the
normals with the Dirichlet components zeroed. \\
\code{UpdateFlux} & Called by \code{Update} and
\code{UpdateWithAcceleration}: gather the velocity,
apply the curl--curl, and call \code{BuildBoundaryPhys}. \\
\code{GatherBoundaryElements} & Perform the indexed copy in
\eqref{hopbc:eq:volume-gather}. \\
\code{CurlCurlOp::Apply} & Evaluate \eqref{hopbc:eq:curlcurl} on boundary-adjacent
volume elements. \\
\code{HOPressureVolumeCombine}\newline\code{Kernel} & Combine the base/predictor
and curl--curl terms on the restricted volume domain, before extraction (see
\S\ref{hopbc:sec:trace-and-normal}). \\
\code{BuildBoundaryPhys} & Combine the flux on the restricted volume, extract
it onto every region, and contract it with the normal; on a region with a
velocity component that is not Dirichlet, extract \(\vect P-\vect R\) too and
subtract its contraction with the normals with the Dirichlet components
zeroed. \\
\code{PhysTraceExtractBndCond}\newline\code{::Apply} & Trace the combined flux to boundary
quadrature points, as in \eqref{hopbc:eq:trace-extract}. \\
\code{HOPressureBoundaryCombine}\newline\code{Kernel} & Contract the traced
flux with \(\vect n\), as in \eqref{hopbc:eq:pointwise-kernel}. \\
\code{SubtractBoundaryVelocity} & Evaluate \(\vect u_D(t)\) on the boundary
and subtract \(s\vect n\cdot\vect u_D\). \\
\code{BuildBoundaryCoeffs} & Apply \eqref{hopbc:eq:boundary-iproduct}, concatenate
regions, and gather into domain-block order. \\
\code{v\_Apply} & Signed and grouped scatter-add in
\eqref{hopbc:eq:weak-scatter}, at the interleave width of the forcing. \\
\bottomrule
\end{tabular}
\end{center}

\section{Execution-space considerations}

The same mathematical sequence is used for Serial, AVX, CUDA, and HIP:
\begin{itemize}
  \item setup maps are built once and made available in the selected memory
        space;
  \item only boundary-adjacent volume elements are copied for curl--curl;
  \item the flux is combined in one kernel before extraction, and contracted
        with the normal in one kernel after it;
  \item work is batched by compatible source, curl--curl, boundary-region, and
        boundary-block layouts;
  \item the restricted fields take the interleave width of the velocity,
        1 or that of the execution space, so that the curl--curl reads and
        writes them without reshaping; the gather copies through index maps
        built on the host for the interleave widths of the source and the
        restricted block, the first time the pair is seen, and kept, and
        numbers its threads by restricted point, the element lane varying
        fastest, so that its writes are coalesced; it
        runs on the stream of each restricted block, which waits on an event
        recorded on the stream of its source block, so the host never waits;
        after that, every block of the restricted volume and of each region
        runs on its own stream;
  \item the extraction orders itself after the streams of the restricted
        blocks it reads, and the streams of the region blocks after itself,
        with events (see \S\ref{sec:phystraceextractbndcond});
  \item the boundary coefficients of the regions are copied into one
        compact array on the streams of the region blocks, and gathered into
        the layout of each domain block on the stream of that block; the
        region and domain block streams wait on one another through the
        operator's own stream before and after the gathers;
  \item the final scatter separates unique coefficient targets from shared
        corner targets, avoiding atomics while retaining the correct sum;
        its target indices are built on the host, for the interleave width
        of a block of the pressure forcing, 1 or that of the execution
        space, the first time a block at that width is seen, so the forcing
        is updated in its own layout, never reshaped.
\end{itemize}
No solver-level backend-specific branch changes equations
\eqref{hopbc:eq:implemented-hop} or \eqref{hopbc:eq:weak-scatter}.

\section{Summary}

The complete discrete operation may be read as
\begin{equation}
\boxed{
\begin{aligned}
\vect C_K
  &=\nabla\times(\nabla\times\vect V_K),\\
g^f_a
  &=s\,\vect n^f_a\cdot
    \left[\mathrm{Extract}_f(\vect P_K-\nu\vect C_K)(x_a)-\vect u_D(x_a,t)\right],\\
\vect b_f
  &=B_f^{\mathsf T}W_fJ_f\vect g_f,\\
\vect r_{\mathrm{pressure}}
  &\mathrel{+}=\sum_f E_f^{\mathsf T}\vect b_f.
\end{aligned}}
\end{equation}
The first line is a volume derivative, the second moves the strong vector flux
to the boundary and forms its normal component, the third converts the
physical boundary flux into a weak load, and the fourth moves that load back
into the volume coefficient system.

\section{Known follow-ups}

The items below do not change \eqref{hopbc:eq:implemented-hop} or which regions
are selected; they concern performance and state.

\begin{itemize}
  \item \textbf{Velocity-expression time-dependence.} Every boundary
  velocity expression is re-evaluated every stage, by
  \code{SubtractBoundaryVelocity}. The Dirichlet and Neumann operators of the
  library, \code{DirBndCondOpImpl} and \code{NeuBndCondOpImpl}, skip it
  for the conditions that are not time-dependent, as recorded by
  \code{CGBndCondMapBuilder}.
  \item \textbf{Velocity history at volume size.} On a rank with a region
  of this condition or of \code{HOutflow}, \code{VelocityCorrectionScheme} keeps
  a history of the velocity as deep as the implicit solves of the time
  integration scheme combine explicit contributions, and the field it is
  extrapolated into,
  on the whole volume, although the operators read only the elements adjacent
  to their regions. Keeping it on those elements instead needs each operator
  to gather the history itself, onto its own restricted elements, and to take
  the extrapolated velocity in that layout. The same holds for the history
  of the explicit forcing kept for a region with a velocity component that
  is not Dirichlet.
\end{itemize}

